\section{Methodology}
In this research we aim to automate parameter-specific temporal resolutions within the Fully Flexible Temporal Resolution (FFTR) framework. The methodology consists of three main components: 
In section \ref{section:Methodology.Hierarchical Time-Series Partitioning} we explain the four different clustering methods to create these flexible temporal resolutions. Then in section \ref{section:Methodology.Validation Methodology} we show statistical measurements (a priori) to evaluate these temporal resolutions. Lastly we show how we verify the introduced error of our methods (regret) in section \ref{section:Methodology.Regret} 

TODO section about hybride RP with my clustering

\subsection{Hierarchical Time-Series Partitioning}
\label{section:Methodology.Hierarchical Time-Series Partitioning}

We apply chronological preserving hierarchical clustering to reduce the temporal resolution of demand and renewable generation profiles (wind and solar), while preserving the statistical and operational characteristics that matter most for energy system optimisation. Unlike conventional approaches that partition the time horizon into uniform blocks of equal width, our method allows variable-width partitions that adapt to the structure of the data.

\subsubsection{Chronology-Preserving hierarchical clustering}

We treat each time step $t$ as an initial singleton cluster and represent the full time horizon as a list of clusters. At each iteration, we merge the adjacent pair $(c_1, c_2)$ that minimizes the Ward dissimilarity


\begin{equation}
\label{eq:ward_distance}
d(c_1, c_2) = \frac{n_1 \, n_2}{n_1 + n_2} \left\| \mu_1 - \mu_2 \right\|^2,
\end{equation}
\noindent
where $n_i$ and $\mu_i$ denote the size and centroid of cluster $c_i$, respectively. Restricting merges to chronologically adjacent clusters, so every resulting partition corresponds to the same order of the original time series. The algorithm continues until the number of clusters reaches a user-specified target $num_clusters$. A nice property of this algorithm is that it is deterministic. It will always terminate and always generates the same output for a specific input. 

There are two ways we can cluster our temporal data. We can either cluster the profiles (e.g.\ demand, wind, and solar) independently of each other. Or we can group these profiles together and cluster them dependently, resulting in the same temporal resolution. Note that when we do this we only group profiles together per location. Doing this per-location profile dependency creates less flexibility and therefore a result with a bigger summed wards dissimilarity. We explain later why this is preferable when doing experiments. 

When computing the clusters with a per-location profiles dependency, Wards desimilarity function becomes multi-dimensional. where for $D$ co-evolving profiles (e.g.\ demand, wind, and solar), we represent each cluster by a $D$-dimensional centroid and compute the dissimilarity in the joint feature space. This means that calculating $\left\| \mu_1 - \mu_2 \right\|^2$ we need to first calculate the difference vector by substracting $\mu_1$ by $\mu_2$. Then we square each ellement in that vector. And finally we sum all values.

After merging two clusters, the centroid $\mu$ of the new cluster is defined as the centroid of all values inside of this cluster.

This algorithm can be efficiently implemented in a linked-list structure abstracted as a priority queue. 

\subsubsection{Extreme-Event Preservation Strategies}

A known limitation of the clustering method described in the previous section is that it tends to absorb extreme values. Since the value of each cluster is the average of all the values within the cluster, the distribution of the values of all clusters decreases the more we cluster. In section \ref{section:Methodology.Validation Methodology} we show this visually. An intuition why this is a bad thing is that the full temporal resolution will have more temesteps with a higher demand then the aggregated version. Therefore when solving the model with the aggrigated version, the model will think that at these extreme timesteps, we need less electricity then we actually do. Therefore we underestimate the electricity demand in these timesteps, where it can be really costly to not have enough power. The same holds for the near-zero renewable generation periods, where the availability will be overestimated.

To combat this systematic bias, our method tries to preserve these extreme events. We define threshold-based extreme flags for each profile dimension:

\begin{itemize}
    \item \textbf{Demand}: a time step is flagged as extreme if its value exceeds 
    the $p_{\text{high}}$ percentile of the full demand profile.
    \item \textbf{Renewables (wind, solar)}: a time step is flagged as extreme if 
    its value falls below the $p_{\text{low}}$ percentile of the respective 
    generation profile.
\end{itemize}

\noindent We compare three strategies for incorporating this extreme-event 
information into the clustering procedure:

\paragraph{Unconstrained Clustering (UC).}
The baseline strategy applies Ward hierarchical clustering with no awareness of extreme 
events. Every adjacent pair of clusters is a candidate for merging, and the 
algorithm selects solely on the basis of Ward dissimilarity. Extreme time steps 
may be freely absorbed into non-extreme partitions, and the representative value 
of every partition is its centroid.

\paragraph{Post-hoc Extreme Correction (PEC).}
This strategy runs the same unconstrained hierarchical clustering but adjusts the 
representative value of each partition after clustering is complete. For demand 
profiles, if the maximum observed value within a partition exceeds the high 
threshold, the partition's representative is set to that maximum rather than the 
centroid. Symmetrically, for renewable profiles, if the minimum observed value 
falls below the low threshold, the minimum replaces the centroid. Formally, the 
representative value for dimension $j$ of cluster $c$ is

\[
\hat{v}_{c,j} = \begin{cases}
\max_{t \in c} x_{t,j} & \text{if } \max_{t \in c} x_{t,j} \geq \tau^{\text{high}}_j 
                          \text{ and profile } j \text{ is demand,} \\
\min_{t \in c} x_{t,j} & \text{if } \min_{t \in c} x_{t,j} \leq \tau^{\text{low}}_j 
                          \text{ and profile } j \text{ is renewable,} \\
\mu_{c,j}              & \text{otherwise.}
\end{cases}
\]

This correction ensures that at least one partition in each high-extreme or 
low-extreme region exposes the true peak or trough to the optimisation model, 
but it does so conservatively: the partition boundaries themselves are unchanged, 
so extreme time steps may still be merged with non-extreme neighbours.

\paragraph{Extreme-Aware Clustering (EAC).}
Our most expressive strategy integrates extreme-event information directly into 
the merge priority. We change the order of which two clusters to merge with a two-level key: each 
candidate merge is first ranked by its \emph{extreme conflict score}---the 
number of profile dimensions for which exactly one of the two clusters contains 
an extreme time step---and then by Ward dissimilarity as a tiebreaker. 
Concretely, for a candidate merge $(c_1, c_2)$:

\[
\text{conflict}(c_1, c_2) = \sum_{j=1}^{d} \mathbf{1}\!\left[ 
    \text{is\_extreme}_{c_1,j} \neq \text{is\_extreme}_{c_2,j} 
\right],
\]

where $\text{is\_extreme}_{c,j}$ is true if any time step in cluster $c$ is 
flagged as extreme in dimension $j$, and $d$ is the number of different profiles we cluster dependently. The algorithm always prefers merges with $\text{conflict} = 0$; only when all remaining candidate merges would combine an extreme partition with a non-extreme one does it accept such a merge, and even then it selects the least disruptive option by Ward dissimilarity. The same post-hoc representative correction as in PEC is then applied to the final partitions.

This design means that extreme time steps aim to merge with 
other extreme time steps before they are ever merged with the surrounding 
background, giving the optimisation model clean, isolated partitions that 
faithfully represent critical operating conditions.

Note that directly modifying wards dissimilarity function to take this representative value $\hat{v}_{c,j}$ instead of $\mu_{cj}$ is not possible. It would break the mathimatical property that makes wards method work; formally stated as:

\[
SSE(c_1 \cup c_2) = SSE(c_1) + SSE(c_2) + d(c_1,c_2)
\]
\noindent
where $SSE(c)$ means Sum Squared Error of cluster $c$ and $d(c_1,c_2)$ is wards desimilarity function described in equation \ref{eq:ward_distance}.

\subsubsection{Summary}

Table~\ref{tab:clustering_strategies} summarises the three strategies. All three 
share the same chronology-preserving Ward hierarchical clustering backbone; they differ 
only in how---and whether---they account for extreme events during and after 
clustering.

\begin{table}[h]
\centering
\caption{Comparison of the three time-series partitioning strategies.}
\label{tab:clustering_strategies}
\begin{tabular}{lccc}
\toprule
\textbf{Strategy} & \textbf{Abbreviation} & 
\makecell{\textbf{Extreme-aware}\\\textbf{merging}} & 
\makecell{\textbf{Extreme representative}\\\textbf{correction}} \\
\midrule
Unconstrained Clustering    & UC  & \texttimes & \texttimes \\
Post-hoc Extreme Correction & PEC & \texttimes & \checkmark \\
Extreme-Aware Clustering    & EAC & \checkmark & \checkmark \\
\bottomrule
\end{tabular}
\end{table}


\subsection{Validation Methodology}
\label{section:Methodology.Validation Methodology}

Aggregation quality is assessed using both statistical and operational validation methods. Statistically, load duration curves are computed from the aggregated time series and compared to those obtained from the full-resolution data. These curves provide insight into how well the aggregation preserves the distributional properties of demand and renewable availability.

Operational validation is measured by relative regret, which follows established best practices in the literature. Investment decisions obtained from the aggregated optimization model are fixed, after which the system is re-optimized or simulated using the original high-resolution time series. Feasibility losses are quantified using Energy Not Served (ENS), allowing direct assessment of whether aggregation-induced simplifications lead to overly optimistic investment outcomes.

\subsection{Regret}
\label{section:Methodology.Regret}

\subsection{Exploration of Peak- and Low-Preserving Clustering}

To improve the representation of extreme events, an alternative clustering variant was explored that explicitly separates peaks and low-demand or low-generation periods. This method modifies the Ward dissimilarity by preventing merges across quantile-based thresholds. Specifically, cluster merges that would combine values across predefined lower and upper quantiles are penalized or forbidden.

The motivation for this approach was to improve the fidelity of load duration curves and extreme value representation. However, empirical results showed that this method led to increased ENS during validation and higher computational effort compared to standard Ward clustering. These findings suggest that rigidly separating extremes can negatively affect the temporal coherence of partitions.

\subsection{Planned Extensions: Hybrid Aggregation and Representative Periods}

Building on these findings, several methodological extensions are planned.

First, hybrid aggregation strategies will be investigated by combining the fully flexible partitions with alternative within-partition aggregation rules. Inspired by the hybrid down-sampling methods proposed by Brafine et al.~\cite{Brafine2024}, this includes selectively using minimum and maximum values for partitions associated with peaks and low extremes, while retaining integral aggregation elsewhere. This approach aims to preserve extreme events without fragmenting the temporal structure excessively.

Second, the interaction between time-series partitioning and representative period selection will be systematically explored. Two aggregation orders will be compared: (i) applying fine-grained partitions before constructing representative periods, and (ii) first clustering representative periods and subsequently refining them using within-period partitioning. While partitioning prior to representative period selection may reduce the uniqueness of periods and degrade clustering quality, applying partitions afterward may distort the internal structure of representative periods. This trade-off has not been systematically evaluated in existing work and will be explicitly addressed in this thesis.

\subsection{maybe also want to concider: (lets discus)}
\subsubsection{Variable Coupling in Partition Construction}

The current implementation constructs partitions independently for demand, wind, and solar time series. While this allows variable-specific resolution control, it may fail to capture joint temporal patterns and correlations between variables. As a further extension, a combined partitioning approach will be explored, in which a single hierarchical clustering is constructed based on a multivariate representation of demand and renewable availability. The performance of independent versus joint partitions will be evaluated in terms of aggregation quality, computational efficiency, and system-level optimization outcomes.


\subsubsection{more case studies}
probably overfitting for dataset \ref{section:experiments:dataset}. should i investigate more scenarios?
